\chapter{Class-Based Two-Button Website}
\ChapterLabel{ClassBasedWebsite}
\index{JavaScript!classes}
\index{UML}

\section{Overview}
The original website attached two anonymous arrow functions directly to the
Blue and Red buttons. Each function repeated the same logic with a different
color. This chapter replaces those inline handlers with three ES2015 classes in
\texttt{public/js/colorButtons.js}, loaded by the page as a JavaScript module.
The visible behavior is unchanged.

\section{Original Design}
\begin{minted}{html}
<script>
  document.getElementById("blueBtn").addEventListener("click", () => {
    document.body.style.backgroundColor = "blue";
  });

  document.getElementById("redBtn").addEventListener("click", () => {
    document.body.style.backgroundColor = "red";
  });
</script>
\end{minted}

\section{Class-Based Design}
\begin{itemize}
\item \textbf{\texttt{BackgroundChanger}} owns the element whose background
changes (the page body) and is the only class that touches its style.
\item \textbf{\texttt{ColorButton}} binds one button element to one color and
delegates the change to a \texttt{BackgroundChanger}.
\item \textbf{\texttt{ColorButtonApp}} holds the button configuration,
creates
one shared \texttt{BackgroundChanger} and one \texttt{ColorButton} per entry,
and binds them. Adding a third color is a one-line change to
\texttt{BUTTONS}.
\end{itemize}

\subsection{UML Class Diagram}
Figure~\ref{Figure:ColorButtonsUML} shows the updated design. A filled diamond
marks composition: the application creates and owns its buttons and its
background changer. The plain arrow is a directed association: each button
holds a reference to the shared changer. An underlined attribute is static;
\texttt{-} marks a field that is private by convention and \texttt{+} a public
method. The same diagram is kept as Mermaid
source in \texttt{docs/uml-class-diagram.md}, which GitHub renders directly.

\begin{figure}[htbp]
\centering
\begin{tikzpicture}[
  umlclass/.style={rectangle split, rectangle split parts=3, draw, thick,
    fill=devopsLightBlue, text width=5.9cm, font=\ttfamily\scriptsize},
  comp/.style={{Diamond[length=3.2mm,width=2.2mm]}-, thick},
  lbl/.style={font=\footnotesize}]
\node[umlclass, anchor=north] (app) at (0,0) {\hfill\textbf{ColorButtonApp}\hfill\null
  \nodepart{second}
  \underline{BUTTONS : Array<\{id, color\}>}\\
  - doc : Document\\
  - changer : BackgroundChanger\\
  - buttons : ColorButton[]
  \nodepart{third}
  + constructor(doc : Document)\\
  + init() : ColorButtonApp};
\node[umlclass, anchor=north] (btn) at (-3.9,-5.0) {\hfill\textbf{ColorButton}\hfill\null
  \nodepart{second}
  - element : HTMLElement\\
  - color : string\\
  - changer : BackgroundChanger
  \nodepart{third}
  + constructor(element, color, changer)\\
  + bind() : void\\
  + handleClick() : void};
\node[umlclass, anchor=north] (chg) at (3.9,-5.0) {\hfill\textbf{BackgroundChanger}\hfill\null
  \nodepart{second}
  - target : HTMLElement
  \nodepart{third}
  + constructor(target : HTMLElement)\\
  + setColor(color : string) : void\\
  + getColor() : string};
\draw[comp] ([xshift=-1.6cm]app.south) -- (btn.north)
  node[lbl, pos=0.85, left] {2} node[lbl, midway, left] {buttons};
\draw[comp] ([xshift=1.6cm]app.south) -- (chg.north)
  node[lbl, pos=0.85, right] {1} node[lbl, midway, right] {changer};
\draw[-{Stealth[length=2.5mm]}, thick] ([yshift=-4mm]btn.text split east) -- ([yshift=-4mm]btn.text split east -| chg.west)
  node[lbl, midway, above] {uses} node[lbl, pos=0.85, below] {1};
\end{tikzpicture}
\caption{UML class diagram of the class-based two-button website.}
\label{Figure:ColorButtonsUML}
\end{figure}

\section{Implementation}
\noindent\textbf{public/js/colorButtons.js} --- the three classes, exported as
an ES module so the browser and the unit tests load the same file:
\inputminted{javascript}{../website/color-buttons-app/public/js/colorButtons.js}

\medskip\noindent\textbf{public/index.html} --- the page now loads the module
and creates the application in two lines:
\inputminted{html}{../website/color-buttons-app/public/index.html}

\medskip\noindent\textbf{server.js} --- unchanged except for a
\texttt{/health} endpoint used by the deployment checks in
\ChapterRef{AWSDockerDeployment}:
\inputminted{javascript}{../website/color-buttons-app/server.js}

\medskip\noindent\textbf{Dockerfile} --- unchanged from the local Docker
exercise. A new \texttt{.dockerignore} keeps \texttt{node\_modules} and the
tests out of the image.
\inputminted{docker}{../website/color-buttons-app/Dockerfile}

\section{Testing}
Because the classes receive the document and elements through their
constructors rather than reading globals, they can be tested in Node.js with
small stand-in objects, without a browser. The tests use Node's built-in
\texttt{node:test} runner, so no extra dependency is needed.
\inputminted{javascript}{../website/color-buttons-app/test/colorButtons.test.js}
Result of \texttt{npm test}:
\inputminted{console}{../evidence/website-unit-tests.txt}
The module was also loaded in a real browser from the running container, where
clicking \emph{Red} then \emph{Blue} set the body background to
\texttt{red} and then \texttt{blue}.
